\documentclass[12pt]{article}
\usepackage[utf8]{inputenc}
\usepackage[T1]{fontenc}
\usepackage[french]{babel}
\usepackage{amsmath}
\usepackage{amssymb}
\usepackage{xcolor}
\usepackage{tikz}
\usepackage{graphicx}
\usetikzlibrary{trees} % Nécessaire pour l'arbre
\begin{document}

\section*{\underline{Dénombrement}}

\subsection*{I) Cardinal d'un ensemble fini}

\subsubsection*{1) Définition}
\fbox{
\begin{minipage}{0.5\textwidth}
Soit $E$ un ensemble fini non vide. Le cardinal de $E$ est le nombre d'éléments contenus dans $E$. \\
On le note : {\color{red} $Card(E)$}
\end{minipage}
}

\paragraph{Exemples :}
\begin{itemize}
    \item $card \{a\} = 1$ ; $card \{a; b\} = 2$
    \item $card D_{14} = 4$ ; $\dots$
\end{itemize}

\subsubsection*{2) Remarques}
\begin{itemize}
    \item $card \emptyset = 0$
    \item $A \subset B \implies card A \leqslant card B$
\end{itemize}

\subsubsection*{3) Résultat}
$A, B$ ensembles finis \\
{\Large $card(A \cup B) = {\color{red} card A + card B - card(A \cap B)}$}

\vspace{2cm}

\section*{\underline{Exercice}}

Lors d'un sondage chez 100 ouvriers dans une usine concernant la lecture de deux Revues $A$ et $B$ ; On a eu les résultats suivants :
\begin{itemize}
    \item 78 lisent la Revue \textcircled{A}
    \item 52 lisent la Revue \textcircled{B}
    \item 40 lisent les deux Revues.
\end{itemize}

Quel est le nombre des ouvriers qui ne lisent aucune des deux Revues ?

% ... dans le corps du texte :
\begin{figure}[h]
    \centering
    \includegraphics[width=0.7\textwidth]{Screenshot from 2026-04-04 22-06-17.png}
    \caption{Diagramme de Venn du cours}
\end{figure}

\begin{tikzpicture}[thick, scale=1.2]
    % 1. L'ensemble E (le contour extérieur)
    % On dessine une grande ellipse pour l'ensemble E
    \draw (0,0) ellipse (3.5cm and 2.2cm);
    \node at (-2.8, 1.8) {\Large $E$};

    % 2. Séparation interne pour créer A et B
    % On dessine les deux arcs de cercle qui se croisent à l'intérieur de E
    % Arc pour la limite de B
    \draw[blue] (0.2, 2.15) to[out=-100, in=100] (0.2, -2.15);
    
    % Arc pour la limite de A
    \draw[red] (-1, 2.15) to[out=-100, in=100] (1, -2.15);

    % 3. Les flèches d'annotation (exactement comme sur ton image)
    
    % Flèche pour A
    \draw[red, <-, >=stealth] (-2.2, -1) -- (-2.8, -1.8) node[below, black] {\Large $A$};
    
    % Flèche pour B
    \draw[blue, <-, >=stealth] (2.5, -1.2) -- (3.2, -1.8) node[below, black] {\Large $B$};
    
    % Flèche pour l'intersection A ∩ B
    \draw[<-, >=stealth] (0.5, -1.5) -- (0.5, -2.2) node[below] {\Large $A \cap B$};

\end{tikzpicture}

\subsection*{II)Principe fondamentale du dénombrement (Principe multiplicatif)}

\noindent
\textcircled{1} \underline{Exemple 1} : On lance dans l'air une pièce de monnaie 3 fois successives. \\
Quel est le nombre total de cas possibles ?

\vspace{1em}

\noindent
\begin{tabular}{|p{6cm}|p{6cm}|}
\hline
\centering \Large \color{red} l'arbre & \centering \Large \color{red} Loi Produit \tabularnewline
\hline
& \\
% --- DÉBUT DE L'ARBRE ---
\begin{tikzpicture}[level distance=1.5cm,
  level 1/.style={sibling distance=2.5cm},
  level 2/.style={sibling distance=1.2cm},
  level 3/.style={sibling distance=0.6cm},
  thick]
  \node {}
    child {node {P}
      child {node {P}
        child {node {P}}
        child {node {F}}
      }
      child {node {F}
        child {node {P}}
        child {node {F}}
      }
    }
    child {node {F}
      child {node {P}
        child {node {P}}
        child {node {F}}
      }
      child {node {F}
        child {node {P}}
        child {node {F}}
      }
    };
\end{tikzpicture} 
& 
% --- DÉBUT DU TABLEAU LOI PRODUIT ---
\vspace{1cm}
\begin{tabular}{|c|c|c|}
\hline
1\textsuperscript{er} lancer & 2\textsuperscript{e} lancer & 3\textsuperscript{e} lancer \\
\hline
{\color{green!60!black}\Huge 2} & {\color{green!60!black}\Huge 2} & {\color{green!60!black}\Huge 2} \\
\hline
\end{tabular}
\newline \centering $\times$ \hspace{0.8cm} $\times$
\vspace{1cm}

Le nombre de cas possibles : \\
{\color{red} \Large $2 \times 2 \times 2 = \underline{\underline{8}}$}
\tabularnewline
\Large On a {\color{red}\underline{8}} cas possibles. & \\
\hline
\end{tabular}

\end{document}